\section{Introduction}\label{sec:intro}

A local correction mechanism can be time independent even when the
events by which it acts are random. The question is whether
uncoordinated local attempts can suppress the syndrome structures that
eventually produce a logical error. In four dimensions the toric code
offers a particularly concrete test: syndrome edges form even-degree
loops, and a conditional plaquette flip can move their upper corners
towards smaller coordinate height.

Pastawski, Clemente and Cirac proposed a dissipative realization of this
four-dimensional Toom mechanism and supported its memory properties
numerically \cite{PCC}. Drechsler, Oreg and Aharonov proved synchronous
and Floquet memory theorems and asked whether a time-independent
Lindbladian, including that proposal, has a nonzero memory threshold
\cite[Sections 1.5.3 and 1.8]{DOA}. Their own physical noise model is
independent single-qubit Pauli noise: Definition 2.25 gives the sum of
one-qubit jump generators, Definition 2.37 the product Pauli layer, and
Definition 7.7 the corresponding Floquet noise assumption
\cite{DOA}. Theorem~\ref{thm:main} therefore answers their
time-independent question in their model.

Theorem~\ref{thm:correlated} goes beyond that model: the exponential
lifetime and fidelity bound also hold under spacetime local stochastic
Pauli noise on unit time cells, allowing bounded-range multi-qubit and
adaptive faults. The noise law must be independent of the hidden
logical input and its reference. For this extension the physical
description is \emph{Lindbladian correction, trajectory-defined noise}.
Our strongest result, Theorem~\ref{thm:general}, proves exponential
storage under any sufficiently weak \emph{additive} finite-range
Markovian Lindblad noise generator, including coherent Hamiltonian
terms and non-unital noise such as amplitude damping.
It controls all six logical qubits and an untouched reference.
The coherent threshold is on the amplitude \(|h|\), not \(h^2\).
Drechsler, Oreg and Aharonov reserve \emph{fully autonomous} memory for
protection against general local noise by a time-independent, local,
translationally invariant Lindbladian
\cite[Section 1.5.1, Table 1]{DOA}.
Theorem~\ref{thm:general} meets that definition for additive finite-range
Markovian noise, with translation invariance for homogeneous noise;
its bound also permits spatially inhomogeneous noise.

In Section 1.4.1 they describe the setting outside thermal equilibrium
as having \textquotedblleft no known rigorous proof of a fully
autonomous quantum memory\textquotedblright\ \cite{DOA}.
The thermal precedent is the four-dimensional toric code under
Davies dynamics studied by Alicki, M.~Horodecki, P.~Horodecki and
R.~Horodecki \cite{AHHH2010}.
Here the added noise generator is required to generate CPTP evolution
and the correction remains exact. Perturbations of the correction's
own jump operators, such as an imperfect syndrome projector, remain
outside the result, as does a general non-Markovian bath.
The former belongs to the broader superoperator-stability question
in Drechsler--Oreg--Aharonov Section 1.8 \cite{DOA}.

The independent-noise conclusion has two related forms. On the finite four-torus a
single, continuously evolving device retains its encoded information
for an exponential time, with an exponentially small worst-case
fidelity error. On the infinite lattice, the syndrome process started
empty has exponential component-diameter tails at any specified
location, uniformly over all physical times. The infinite statement
concerns syndrome geometry; it does not assign a logical qubit to the
infinite lattice.

\subsection{Relation to existing work}

The directed geometry and the contour method originate with Toom
\cite{Toom}. Bennett and Grinstein studied the stabilizing role of
irreversibility in local stochastic dynamics \cite{BennettGrinstein}.
Bramson and Gray developed renormalization comparisons with growth and
erosion processes \cite{BramsonGray}; Gray proved continuous-time
stability results for attractive spin systems
\cite{Gray}. These works establish a substantial classical
background for asynchronous stability.

The present syndrome process is not an attractive spin system under
edgewise inclusion. An enabled plaquette operation removes the two
tested edges and toggles the opposite two; an opposite edge can
disappear if it was already marked. Thus inclusion of edge sets does
not give the order needed for a direct attractive-process comparison.
We instead construct an auxiliary monotone process on spatial blocks
and verify its implications for every admissible starting syndrome in
a slab.

Chakraborty and Lee use the same kind of slab comparison to study
finite-range monotone continuous-time dynamics
\cite[Theorem 1 and Appendix B]{ChakrabortyLee}. The local comparison here
is established directly from the parity of syndrome vertices and the
waiting times of the physical Poisson clocks. Swart, Szab\'o and
Toninelli provide the deterministic contour construction and its
enumeration \cite[Theorem 22 and Lemmas 38--40]{SST}. We apply that
construction to the block process, with a separate calculation for the
finite-range dependence of the bad-bin field.

For quantum codes, Dennis, Kitaev, Landahl and Preskill developed the
homological account of topological quantum memory and discussed local
correction in four dimensions \cite{DKLP}. From Ahn's thesis we use only
the five linear forms in Eq.~(7.17) \cite{Ahn}; their required
inequalities are checked directly in Section~\ref{sec:contours}.
No dynamical or threshold estimate from that thesis is an input to
this proof. Kubica and Preskill proved threshold results for
synchronous cellular-automaton sweep decoders \cite{KubicaPreskill}.
Ponarovsky treats local stochastic Pauli noise with synchronous sweeps
and noiseless bounded-depth classical processing
\cite[Sections 1.1--1.3, Theorems 6.3 and 7.3]{Ponarovsky}.
That paper also notes gaps in Ahn's four-dimensional Toom proof and
points to a companion discussion \cite[Section 1.1]{Ponarovsky}.
Its timing and processing resources differ from independent Poisson
correction attempts.

Kumar and Weimer study a two-dimensional dissipative toric code with
a dynamical classical field \cite{KumarWeimer}. D\"unnweber, Styliaris
and Trivedi give a two-dimensional autonomous cellular-automaton
construction and outline a continuous-time realization \cite{DST}.
These are related constructions, not inputs to the estimates below.
Joint charging of fault tuples is standard in local-stochastic
fault tolerance; the correlated extension combines it with the
asynchronous clearance, contour and overlap arguments.

We credit the quantum dissipative proposal and its numerical evidence
to Pastawski, Clemente and Cirac \cite{PCC}, and the synchronous and
Floquet theorems and the explicit time-independent question to
Drechsler, Oreg and Aharonov \cite{DOA}. The method used here is standard
in kind: a block comparison, Toom contours, a Peierls estimate and a
homological decoding argument. Its application to the non-monotone
four-dimensional toric-code syndrome dynamics appears new.

\subsection{What the probability estimate controls}

We prove, in either CSS sector, the history bound
\begin{equation}\label{eq:intro-history}
 \begin{aligned}
 &\Prb\{\text{a syndrome component of diameter }\ge L/100
             \text{ occurs by time }t\}\\
 &\hspace{3em}\le C_0(1+t)e^{-\beta L}
 \qquad(t\ge0).
 \end{aligned}
\end{equation}
The constants and the lower bound on \(L\) are uniform over all allowed
spatial choices of the noise rates. The inequality is understood with
its right side replaced by its minimum with one. It is a statement
about all actual jumps up to a deterministic horizon, rather than only
the syndrome at regularly spaced observations.

A fixed filling of each small syndrome component defines a decoded
homology class. A single plaquette flip cannot change this class when
the components immediately before and after the flip are small.
Equation~\eqref{eq:intro-history} therefore controls the entire decoded
history. Applying this argument in both sectors gives an event on
which every code state, including a state entangled with a reference,
is recovered exactly. The worst-case fidelity bound follows from the
probability of that event.

\subsection{Why overlapping comparisons are needed}

A short backward dependency cone on a large torus lifts injectively,
date by date, to the integer lattice. A contour in such a cone has
distinct physical sinks, so each sink can be charged to a rare
collection of Poisson increments. Over a much longer cone, distinct
lifted sinks might be copies of the same physical bin. Multiplying
their costs would then be invalid.

We avoid this repetition by estimating two comparisons initialized one
epoch apart. They use the same maps during their overlap. If they
still disagree at its end, a path of zeros of the older comparison
crosses an entire epoch in time. Localized spacetime witnesses make
such a path exponentially unlikely while the older comparison is
still young enough to lift. Agreement can then be propagated from one
epoch to the next. The comparisons are functions of the marks; they do
not restart the physical memory.

There is a genuine reason to retain the clean-initialization
hypothesis. Some wrapping, even-degree syndromes are frozen under
every noiseless correction attempt. Consequently the argument does
not assert uniform contraction from arbitrary torus syndromes.
Section~\ref{sec:discussion} gives an explicit example.

\subsection{Organization and conventions}

Section~\ref{sec:model} specifies the generator, fidelity convention and
statements. Section~\ref{sec:local} proves the asynchronous clearance
and the uniform local mark test. Section~\ref{sec:contours} constructs
localized certificates and proves the centered path estimate.
Section~\ref{sec:windows} establishes
\eqref{eq:intro-history} by overlapping windows.
Sections~\ref{sec:decoder} and \ref{sec:quantum} construct the readout and
complete the quantum theorem. Section~\ref{sec:infinite} gives the
infinite-volume construction and its uniform estimates.
Section~\ref{sec:correlated} states and proves the correlated-noise
extension, with its time-grid dependence and joint witness costs.
Section~\ref{sec:general} proves the strongest, additive Markovian-noise
theorem by retaining coherent cross terms, resumming optional faults,
and counting dense unions of witnesses.

Chains and boundaries have coefficients in \(\F\). A syndrome component
means a component of marked edges joined at common endpoints.
Diameters are ambient sup-metric diameters of their vertex sets, using
periodic coordinate distances on a torus. Constants denoted by \(C\)
may be enlarged from one inequality to the next. Every fixed block
parameter is chosen independently of \(L\).
